\textbf{Phần 1. Trắc nghiệm lựa chọn}
\setcounterpageref{ex}{0}

\begin{ex} %[1D7V25]
Bút chì màu đỏ dài 9 cm, bút chì màu xanh dài 7 cm. Hỏi bút chì màu đỏ dài hơn bút chì màu xanh bao nhiêu xăng-ti-mét?
\choice
{$1\text{ cm}$}
{$3\text{ cm}$}
{\True $2\text{ cm}$}
{$4\text{ cm}$}
\loigiai{
Bút chì đỏ dài hơn bút chì xanh là: $9 - 7 = 2$ (cm).
\textbf{Đáp số:} $2$.\par
\textbf{Đáp án của câu hỏi là C}
}
\end{ex}

\begin{ex} %[1D7V25]
Băng giấy màu vàng dài 10 cm, băng giấy màu tím dài 6 cm. Hỏi băng giấy màu tím ngắn hơn băng giấy màu vàng bao nhiêu xăng-ti-mét?
\choice
{\True $4\text{ cm}$}
{$3\text{ cm}$}
{$5\text{ cm}$}
{$2\text{ cm}$}
\loigiai{
Băng giấy màu tím ngắn hơn là: $10 - 6 = 4$ (cm).
\textbf{Đáp số:} $4$.\par
\textbf{Đáp án của câu hỏi là A}
}
\end{ex}

\begin{ex} %[1D7V25]
Đoạn thẳng AB dài 8 cm, đoạn thẳng CD dài hơn đoạn thẳng AB là 2 cm. Hỏi đoạn thẳng CD dài bao nhiêu xăng-ti-mét?
\choice
{$0\text{ cm}$}
{$20\text{ cm}$}
{$9\text{ cm}$}
{\True $10\text{ cm}$}
\loigiai{
Độ dài đoạn thẳng CD là: $8 + 2 = 10$ (cm).
\textbf{Đáp số:} $10$.\par
\textbf{Đáp án của câu hỏi là D}
}
\end{ex}

\begin{ex} %[1D7V25]
Cây thước kẻ dài 20 cm, que tính dài 10 cm. Hỏi cây thước kẻ dài hơn que tính bao nhiêu xăng-ti-mét?
\choice
{$0\text{ cm}$}
{\True $10\text{ cm}$}
{$20\text{ cm}$}
{$9\text{ cm}$}
\loigiai{
Thước kẻ dài hơn que tính là: $20 - 10 = 10$ (cm).
\textbf{Đáp số:} $10$.\par
\textbf{Đáp án của câu hỏi là B}
}
\end{ex}

\begin{ex} %[1D7V25]
Sợi dây thứ nhất dài 15 cm, sợi dây thứ hai ngắn hơn sợi dây thứ nhất 5 cm. Hỏi sợi dây thứ hai dài bao nhiêu xăng-ti-mét?
\choice
{$0\text{ cm}$}
{$20\text{ cm}$}
{\True $10\text{ cm}$}
{$9\text{ cm}$}
\loigiai{
Độ dài sợi dây thứ hai là: $15 - 5 = 10$ (cm).
\textbf{Đáp số:} $10$.\par
\textbf{Đáp án của câu hỏi là C}
}
\end{ex}

\begin{ex} %[1D7V25]
Chiếc đũa dài 18 cm, chiếc thìa dài 12 cm. Hỏi chiếc thìa ngắn hơn chiếc đũa bao nhiêu xăng-ti-mét?
\choice
{\True $6\text{ cm}$}
{$5\text{ cm}$}
{$7\text{ cm}$}
{$4\text{ cm}$}
\loigiai{
Chiếc thìa ngắn hơn chiếc đũa là: $18 - 12 = 6$ (cm).
\textbf{Đáp số:} $6$.\par
\textbf{Đáp án của câu hỏi là A}
}
\end{ex}

\begin{ex} %[1D7V25]
Đoạn thẳng MN dài 7 cm, đoạn thẳng PQ dài 7 cm. Hỏi đoạn thẳng MN dài hơn đoạn thẳng PQ bao nhiêu xăng-ti-mét?
\choice
{$1\text{ cm}$}
{$2\text{ cm}$}
{$3\text{ cm}$}
{\True $0\text{ cm}$}
\loigiai{
Vì hai đoạn thẳng có độ dài bằng nhau nên đoạn thẳng MN dài hơn đoạn thẳng PQ là $0$ cm.
\textbf{Đáp số:} $0$.\par
\textbf{Đáp án của câu hỏi là D}
}
\end{ex}

\begin{ex} %[1D7V25]
Gang tay của mẹ dài 18 cm, gang tay của bé dài 11 cm. Hỏi gang tay của mẹ dài hơn gang tay của bé bao nhiêu xăng-ti-mét?
\choice
{$6\text{ cm}$}
{\True $7\text{ cm}$}
{$8\text{ cm}$}
{$5\text{ cm}$}
\loigiai{
Gang tay của mẹ dài hơn của bé là: $18 - 11 = 7$ (cm).
\textbf{Đáp số:} $7$.\par
\textbf{Đáp án của câu hỏi là B}
}
\end{ex}

\begin{ex} %[1D7V25]
Bút mực dài 14 cm, tẩy chì dài 4 cm. Hỏi bút mực dài hơn tẩy chì bao nhiêu xăng-ti-mét?
\choice
{$0\text{ cm}$}
{$20\text{ cm}$}
{\True $10\text{ cm}$}
{$9\text{ cm}$}
\loigiai{
Bút mực dài hơn tẩy chì là: $14 - 4 = 10$ (cm).
\textbf{Đáp số:} $10$.\par
\textbf{Đáp án của câu hỏi là C}
}
\end{ex}

\begin{ex} %[1D7V25]
Một thanh gỗ dài 19 cm, bác thợ cưa bớt đi một đoạn còn lại 10 cm. Hỏi đoạn gỗ bị cưa bớt dài bao nhiêu xăng-ti-mét?
\choice
{\True $9\text{ cm}$}
{$8\text{ cm}$}
{$10\text{ cm}$}
{$7\text{ cm}$}
\loigiai{
Đoạn gỗ bị cưa bớt dài là: $19 - 10 = 9$ (cm).
\textbf{Đáp số:} $9$.\par
\textbf{Đáp án của câu hỏi là A}
}
\end{ex}

\begin{ex} %[1D7V25]
Đoạn thẳng thứ nhất dài 6 cm, đoạn thẳng thứ hai dài 8 cm, đoạn thẳng thứ ba dài 10 cm. Hỏi đoạn thẳng dài nhất dài hơn đoạn thẳng ngắn nhất bao nhiêu xăng-ti-mét?
\choice
{$3\text{ cm}$}
{$5\text{ cm}$}
{$2\text{ cm}$}
{\True $4\text{ cm}$}
\loigiai{
Đoạn dài nhất là 10 cm, ngắn nhất là 6 cm: $10 - 6 = 4$ (cm).
\textbf{Đáp số:} $4$.\par
\textbf{Đáp án của câu hỏi là D}
}
\end{ex}

\begin{ex} %[1D7V25]
Một cây mía dài 10 dm, bố chặt bớt đi 3 dm. Hỏi cây mía còn lại dài bao nhiêu đề-xi-mét?
\choice
{$6$}
{\True $7$}
{$8$}
{$5$}
\loigiai{
Chiều dài cây mía còn lại là: $10 - 3 = 7$ (dm).
\textbf{Đáp số:} $7$.\par
\textbf{Đáp án của câu hỏi là B}
}
\end{ex}

\textbf{Phần 2. Trắc nghiệm đúng sai}
\setcounterpageref{ex}{0}

\begin{ex} %[1D7V25]
Xét tính đúng hoặc sai của các khẳng định toán học sau đây:
\choiceTF
{\True Sợi len đỏ dài 12 cm, sợi len xanh dài 16 cm. phải cắt bớt sợi len xanh $4\text{ cm}$ xăng-ti-mét để hai sợi len dài bằng nhau.}
{Băng giấy màu xanh dài 5 cm, nối thêm vào một băng giấy màu đỏ dài 4 cm. băng giấy mới dài $8\text{ cm}$ xăng-ti-mét.}
{\True Một cái cột cao 8 gang tay, một cái cây cao 10 gang tay. cái cây cao hơn cái cột $2$ gang tay.}
{Kết quả của bài toán «Con đường từ nhà An đến trường dài 9 bước chân, con đường từ nhà Bình đến trường dài 6 bước chân. Con đường nhà An dài hơn con đường nhà Bình bao nhiêu bước chân» là $2$.}
\loigiai{
\begin{itemize}
    \item Ý a) Đúng: Sợi len xanh dài hơn sợi len đỏ là: $16 - 12 = 4$ (cm). Vậy cần cắt bớt $4$ cm.
\textbf{Đáp số:} $4$.
    \item Ý b) Sai: Kết quả đúng là $9$, không phải $8$. Chiều dài băng giấy mới sau khi nối là: $5 + 4 = 9$ (cm).
\textbf{Đáp số:} $9$.
    \item Ý c) Đúng: Cái cây cao hơn cái cột là: $10 - 8 = 2$ (gang tay).
\textbf{Đáp số:} $2$.
    \item Ý d) Sai: Kết quả đúng là $3$, không phải $2$. Con đường nhà An dài hơn là: $9 - 6 = 3$ (bước chân).
\textbf{Đáp số:} $3$.
\end{itemize}
\begin{itemize}
    \item a) \textbf{Đúng}.
    \item b) \textbf{Sai}.
    \item c) \textbf{Đúng}.
    \item d) \textbf{Sai}.
\end{itemize}
}
\end{ex}

\begin{ex} %[1D7V25]
Xét tính đúng hoặc sai của các bài toán và phép tính sau đây:
\choiceTF
{Đoạn dây nhôm dài 17 cm, Nam uốn thành hình vuông hết 12 cm. đoạn dây nhôm còn thừa dài $4\text{ cm}$ xăng-ti-mét.}
{\True Kết quả của bài toán «Quyển sách Toán dài 24 cm, quyển vở bài tập dài 20 cm. Quyển sách Toán dài hơn quyển vở bao nhiêu xăng-ti-mét» là $4\text{ cm}$.}
{Một cành cây dài 15 cm bị gãy bớt một đoạn dài 5 cm. cành cây còn lại dài $11\text{ cm}$ xăng-ti-mét.}
{\True Một cái cột cao 8 gang tay, một cái cây cao 10 gang tay. cái cây cao hơn cái cột $2$ gang tay.}
\loigiai{
\begin{itemize}
    \item Ý a) Sai: Kết quả đúng là $5$, không phải $4$. Đoạn dây nhôm còn thừa là: $17 - 12 = 5$ (cm).
\textbf{Đáp số:} $5$.
    \item Ý b) Đúng: Quyển sách Toán dài hơn là: $24 - 20 = 4$ (cm).
\textbf{Đáp số:} $4$.
    \item Ý c) Sai: Kết quả đúng là $10$, không phải $11$. Cành cây còn lại dài: $15 - 5 = 10$ (cm).
\textbf{Đáp số:} $10$.
    \item Ý d) Đúng: Cái cây cao hơn cái cột là: $10 - 8 = 2$ (gang tay).
\textbf{Đáp số:} $2$.
\end{itemize}
\begin{itemize}
    \item a) \textbf{Sai}.
    \item b) \textbf{Đúng}.
    \item c) \textbf{Sai}.
    \item d) \textbf{Đúng}.
\end{itemize}
}
\end{ex}

\textbf{Phần 3. Trả lời ngắn}
\setcounterpageref{ex}{0}

\begin{ex} %[1D7V25]
Hộp bút dài 18 cm, chiếc bút chì dài 15 cm. Hỏi hộp bút dài hơn chiếc bút chì bao nhiêu xăng-ti-mét?
\shortans{3}
\loigiai{
\textbf{Đáp số:} $3$\par
Hộp bút dài hơn chiếc bút chì là: $18 - 15 = 3$ (cm).
\textbf{Đáp số:} $3$.
}
\end{ex}

\begin{ex} %[1D7V25]
Một con ốc sên bò lên cây cột cao 10 cm. Ban ngày bò lên được 5 cm, ban đêm bò tiếp lên được 4 cm. Hỏi ốc sên đã bò được tất cả bao nhiêu xăng-ti-mét?
\shortans{9}
\loigiai{
\textbf{Đáp số:} $9$\par
Tổng quãng đường ốc sên đã bò là: $5 + 4 = 9$ (cm).
\textbf{Đáp số:} $9$.
}
\end{ex}

\begin{ex} %[1D7V25]
Chú ốc sên còn cách đỉnh cột bao nhiêu xăng-ti-mét?
\shortans{1}
\loigiai{
\textbf{Đáp số:} $1$\par
Khoảng cách còn lại đến đỉnh cột là: $10 - 9 = 1$ (cm).
\textbf{Đáp số:} $1$.
}
\end{ex}

\begin{ex} %[1D7V25]
Đoạn thẳng AB dài 11 cm, đoạn thẳng CD ngắn hơn AB là 1 cm. Hỏi đoạn thẳng CD dài bao nhiêu xăng-ti-mét?
\shortans{10}
\loigiai{
\textbf{Đáp số:} $10$\par
Độ dài đoạn thẳng CD là: $11 - 1 = 10$ (cm).
\textbf{Đáp số:} $10$.
}
\end{ex}

\begin{ex} %[1D7V25]
Hai thanh gỗ cùng dài 8 cm ghép nối tiếp nhau. Hỏi thanh gỗ mới dài bao nhiêu xăng-ti-mét?
\shortans{16}
\loigiai{
\textbf{Đáp số:} $16$\par
Chiều dài thanh gỗ mới là: $8 + 8 = 16$ (cm).
\textbf{Đáp số:} $16$.
}
\end{ex}

\begin{ex} %[1D7V25]
Chiếc thìa dài 14 cm, chiếc đĩa ăn cơm có đường kính 18 cm. Đường kính chiếc đĩa dài hơn chiếc thìa bao nhiêu xăng-ti-mét?
\shortans{4}
\loigiai{
\textbf{Đáp số:} $4$\par
Đường kính chiếc đĩa dài hơn chiếc thìa là: $18 - 14 = 4$ (cm).
\textbf{Đáp số:} $4$.
}
\end{ex}
